Can a deterministic multitape Turing machine transpose a $d\times d$ matrix with $\lceil\log_2 d\rceil$-bit entries in time
linear in the input length? The question is open, and a negative answer would give the first superlinear lower bound for
off-line integer multiplication on such machines. We ask what a linear-time transposer would need, and we find three
ingredients: curved head motion; depth-robust dependencies, or extreme coding on shallow ones; and storage codes
that beat rate~$\tfrac12$.
We prove that no machine, oblivious or adaptive, whose work heads are all home-based with legs made of boundedly
many constant-rate runs transposes in linear time, and that the macroscopic head motion of every linear-time transposer
contains a linear amount of curved, not countably-piecewise affine, travel (Theorems~C and~Cq). We show that two statements together would
rule out oblivious linear-time transposers: depth reduction of the time--space dependency DAG by a cheap genie, and an
index-coding lower bound on the remaining shallow cones (Theorem~D). We prove depth reduction for every coupling of
home-based walks (Theorem~B), show that packing certificates cannot establish depth-robustness, and give an exact
accounting of label costs on the sharpest open instance, the deep-nesting pair (Lemma~Zq, Theorems~E and~E$_1$). For the index-coding side we study the dual notion, storage codes, on sparse graphs. Rate bounds pull back along
homomorphisms (Theorem~H), so every storage code, nonlinear included, on a graph that maps to $C_5$, to $K_{14/5}$ or to
the Petersen graph has rate at most $\tfrac12$, whatever its degree. The Clebsch graph has guessing number exactly~$10$, so bounded degree alone does not force rate~$\tfrac12$.
On triangle-free graphs of maximum degree~$3$, no linear storage code over any field beats rate~$\tfrac12$ (Theorem~S).
On locally bipartite graphs of bounded degree, linear codes with bounded complex coefficients have rate at most
$\tfrac12+o(1)$, so gains of finite-field linear codes with a fixed number of symbols per vertex need
$p$-torsion (Theorem~A). The entropy-based results on storage codes (Theorem~H, symmetrization, the good targets and the Clebsch graph),
Theorem~S and the exact label accounting (Lemma~Zq, the exact cases of the chain and transfer identities, and
Theorem~E) are formalized in Lean~4, apart from the classical Tutte--Berge formula
used in the matching form of Theorem~S; all other results are proved in full but not formalized. We end with a roadmap: positive answers to depth reduction for interacting walks and to the index-coding
question on shallow cones, which is open for finite-field and nonlinear codes even in its undirected form, would
rule out oblivious linear-time transposers,
and our restricted results say what a construction would have to combine.
